% Section 19.1, translated from sv/chapters/19/exam.tex.

\section{Before the exam}
\label{c19:sec:infor-tentamen}
In the exam you show that you have achieved the course's learning outcomes. This section describes
how the exam works, what it contains and how to prepare for it.

\subsection{Format}
\label{c19:sec:format}
The exam is a written exam, sat in an exam hall. Every question requires a solution and a
justification, with a stated answer, including the unit where relevant. In other words, you must
explain how you arrived at your results.

\subsection{Permitted aids}
\label{c19:sec:hjalpmedel}
\begin{itemize}
  \item Writing materials and a calculator of any kind.
  \item A formula sheet on one A4 sheet (both sides), handwritten or typed.
  \item One handwritten A4 sheet (both sides) with notes of your own choice.
\end{itemize}

\begin{aside}
\note[Note] Mobile phones may not be used during the exam and must be left in the place
indicated.
\end{aside}

\subsection{Points and grade thresholds}
\label{c19:sec:betyg}
The exam gives $25$ points in total. On top of that, up to $1.0$ bonus point can be earned from
the course's two tests: $0.5$ points per test for anyone who scores at least $3.5$ of the test's
$5$ points. The course uses the Swedish grades IG (fail), G (pass) and VG (pass with
distinction), with the thresholds in the table below; a pass grade also requires that all the
course's learning outcomes are achieved.

\begin{narrowtable}{@{}ll@{}}
  \thead{Grade} & \thead{Threshold} \\
  \midrule
  IG & $< 10.0$ points \\
  G & $\geq 10.0$ points \\
  VG & $\geq 17.5$ points \\
\end{narrowtable}

\subsection{Content}
\label{c19:sec:innehall}
The exam is divided into two parts:
\begin{itemize}
  \item \textbf{Part I}: arithmetic, algebra, equations, trigonometry and logarithms
    (Chapters~\ref{c1:ch}--\ref{c10:ch}).
  \item \textbf{Part II}: derivatives, integrals and complex numbers
    (Chapters~\ref{c12:ch}--\ref{c17:ch}).
\end{itemize}

\subsection{Preparation}
\label{c19:sec:forberedelse}
\begin{itemize}
  \item Sit the practice exam in \secref{c19:sec:ovningstentamen} on your own.
  \item Then compare with the answer key and with the worked solutions, in Swedish, published in
    the course repository in \file{exam/practice_exam_solutions.md}.
  \item Review the central parts of the course, especially the areas where you went wrong in the
    practice exam.
  \item Prepare the formula sheet and the sheet of notes (\secref{c19:sec:hjalpmedel}) in good
    time.
\end{itemize}
